\originalsection{场的局部意义与表示}\label{tb:fields}
隔着空间发生相互作用时, 场的描述分成两件事: 源改变周围空间的场, 受作用物体对它所在处的场作出响应. 静态库仑式把这一关系写得像瞬时作用, 但时变扰动的传播需要时间; 第\ref{ch:maxwell}章将从场方程求出传播速度. 这里不把真空场解释成一种必需的机械以太. 场能储能和传递动量, 是可通过受力和能量交换检验的物理对象.

标量场$T(\rr,t)$在每个位置、时刻给一个带单位的数; 温度和电势属于此类. 矢量场$\vct v(\rr,t)$给有大小和方向的量, 其三个分量各自是标量场. 给某个粒子的速度与给整个流体空间的速度场不同: 前者随粒子走, 后者在选定位置观察经过那里的流体.

\begin{example}\label{ex:planet}
原第1章例1.1用半球始终朝向恒星的行星说明标量场. 半径为$R$, 球坐标的$\theta$从北轴量起, $\phi=\pi/2$指向朝阳半球中心, 给定
\[
T(r,\theta,\phi)=\bigl[T_0+T_1\sin^2\theta+T_2(1+\sin\phi)\bigr]e^{-\alpha(r-R)}.
\]
解释各项, 并说明模型的限制.
\end{example}
\begin{solution}
$T_0,T_1,T_2$有温度单位, $\alpha$有长度倒数单位, 因而指数无量纲. 当$T_1,T_2,\alpha>0$, 在表面赤道朝阳中心值为$T_0+T_1+2T_2$, 赤道背阳中心为$T_0+T_1$; 增高$1/\alpha$使温度乘$e^{-1}$. 在固定$\phi$的经线上接近极点, $\sin^2\theta$项下降. 这些结论来自指定函数, 不是大气传热方程的解.

还有一个坐标条件不能省掉: 极点的$\phi$不唯一, 而原式的$T_2\sin\phi$没有随$\sin\theta$消失, 所以它不能直接定义包含极点的单值连续全球温度场. 可把原模型用于离开极点的区域; 若要建立全球平滑示意模型, 可把这一角项改为$T_2(1+\sin\theta\sin\phi)$. 后者是本稿提出的修正, 不是原件公式.
\end{solution}

把标量场画成图时, 等值线只说明同一函数值; 色图用颜色表示大小; 高度图把函数值当成额外绘图高度. 对$V(x,y,0)$的高度图, 竖轴是电势而不是空间中的$z$坐标. 三维标量场还需说明展示的是切面、等值面还是体数据. 比如$1/r$在$r=0$有奇点, 绘图删掉中心小区域只能避开无穷值, 不表示原函数在中心有限.

矢量箭头图的方向可直接读出, 长度是否正比场强要核图例. 场线是处处与矢量场相切的曲线: 若用参数$s$, 可写$\dd\rr/\dd s=\vct F(\rr)$, 或只要求两者平行. 草籽、铁屑和相关纹理能显示局部取向, 却未必直接给出箭头或绝对强度. 尤其铁屑是被磁化的小物体, 展示的是取向, 不是磁场里实际流动的物质.

对多变量场,$\partial_xf$表示保持$y,z,t$固定时沿$x$方向的变化率. $\nabla$是把三个方向微分组合的记号, 散度$\partial_xv_x+\partial_yv_y+\partial_zv_z$是局部有向通量密度, 旋度的$z$分量$\partial_xv_y-\partial_yv_x$是局部小环量密度; 下文将用小盒和小矩形实际导出. 柱坐标$\phi$是绕轴的方位角, $\uv\phi$沿增大角度的圆周切向, 从$+z$看为逆时针. 第\ref{ch:math}章再给完整坐标和积分元素, 此处先知道每一个符号实际测什么.

动画中的纹理移动常用来编码场能的流向, 它不是场线物质沿空间搬运. 粒子的轨迹由力和惯性决定, 场线只显示某时刻场方向; 第\ref{ch:maxwell}章的Poynting矢量才给实际能流量.

\originalsection{源、汇与环流}\label{tb:fieldstructure}
水从喷口流出, 用流体速度场容易想象“源”; 排水口是“汇”; 绕中心旋转是“环流”. 在一个边长很小的直角盒子中, 相对$x$面的向外通量之差为$[v_x(x+\Delta x)-v_x(x)]\Delta y\Delta z$. 加上另两对面并除以体积, 极限是
\[
\nabla\cdot\vct v=\partial_xv_x+\partial_yv_y+\partial_zv_z.
\]
这就是局部净外流率密度. 对光滑恒密度不可压流, 它为零. 若密度会变, 真正的质量守恒是$\partial_t\rho_{\rm f}+\nabla\cdot(\rho_{\rm f}\vct v)=0$, 不能把“速度散度不零”直接说成物质被创造.

类似地, 沿$xy$小矩形逆时针积分, 四条边相消后的首阶项为$(\partial_xv_y-\partial_yv_x)\Delta x\Delta y$. 它是旋度的$z$分量. 另两种矩形给
\[
\nabla\times\vct v=(\partial_yv_z-\partial_zv_y)\uv x
 +(\partial_zv_x-\partial_xv_z)\uv y
 +(\partial_xv_y-\partial_yv_x)\uv z.
\]
“绕着某物弯曲”与局部旋度不零不是同一条件. 例如$\uv\phi/s$在轴外有环状场线, 局部旋度却为零; 绕轴的环路围住被排除的奇点. 是否能由局部无旋推出所有闭环积分零, 还依赖区域与奇点.

\begin{example}\label{ex:fieldplots}
原1.10.1的四种场为$(3,-5)$、$\rr$、$\uv r/r^2$及平面场$(3xy,2y^2-x^2)/r^5$, $r=\sqrt{x^2+y^2}$. 说明其图形, 并给出一组源与环流的可计算比较.
\end{example}
\begin{solution}
第一种处处同向同长; 第二种径向朝外且长度随$r$增大; 第三种径向朝外而大小随$r^{-2}$减少, 原点不定义. 最后一种是沿$+y$的偶极远场形状, 来自$-\nabla(y/r^3)$, 在正、负$y$轴都指向$+y$, 在负半轴为朝向原点, 而在$x$轴朝$-y$, 轴外还有横向分量. 场的方向随位置变, 不能只画一根偶极轴来代替整个场.

\begin{center}\begin{tikzpicture}[>=Stealth,scale=.63]
\begin{scope}[shift={(0,0)}]
\draw[->,gray] (-2.3,0)--(2.4,0) node[right] {$x$};\draw[->,gray] (0,-2.3)--(0,2.4) node[above] {$y$};
\node at (0,-2.9) {$(3,-5)$};
\draw[->,structurecolor] (-2,-2)--(-1.75845,-2.40258);
\draw[->,structurecolor] (-2,-1)--(-1.75845,-1.40258);
\draw[->,structurecolor] (-2,0)--(-1.75845,-0.40258);
\draw[->,structurecolor] (-2,1)--(-1.75845,0.59742);
\draw[->,structurecolor] (-2,2)--(-1.75845,1.59742);
\draw[->,structurecolor] (-1,-2)--(-0.75845,-2.40258);
\draw[->,structurecolor] (-1,-1)--(-0.75845,-1.40258);
\draw[->,structurecolor] (-1,0)--(-0.75845,-0.40258);
\draw[->,structurecolor] (-1,1)--(-0.75845,0.59742);
\draw[->,structurecolor] (-1,2)--(-0.75845,1.59742);
\draw[->,structurecolor] (0,-2)--(0.24155,-2.40258);
\draw[->,structurecolor] (0,-1)--(0.24155,-1.40258);
\draw[->,structurecolor] (0,0)--(0.24155,-0.40258);
\draw[->,structurecolor] (0,1)--(0.24155,0.59742);
\draw[->,structurecolor] (0,2)--(0.24155,1.59742);
\draw[->,structurecolor] (1,-2)--(1.24155,-2.40258);
\draw[->,structurecolor] (1,-1)--(1.24155,-1.40258);
\draw[->,structurecolor] (1,0)--(1.24155,-0.40258);
\draw[->,structurecolor] (1,1)--(1.24155,0.59742);
\draw[->,structurecolor] (1,2)--(1.24155,1.59742);
\draw[->,structurecolor] (2,-2)--(2.24155,-2.40258);
\draw[->,structurecolor] (2,-1)--(2.24155,-1.40258);
\draw[->,structurecolor] (2,0)--(2.24155,-0.40258);
\draw[->,structurecolor] (2,1)--(2.24155,0.59742);
\draw[->,structurecolor] (2,2)--(2.24155,1.59742);
\end{scope}
\begin{scope}[shift={(6,0)}]
\draw[->,gray] (-2.3,0)--(2.4,0) node[right] {$x$};\draw[->,gray] (0,-2.3)--(0,2.4) node[above] {$y$};
\node at (0,-2.9) {$\rr$};
\draw[->,structurecolor] (-2,-2)--(-2.28732,-2.28732);
\draw[->,structurecolor] (-2,-1)--(-2.33992,-1.16996);
\draw[->,structurecolor] (-2,0)--(-2.36667,0.00000);
\draw[->,structurecolor] (-2,1)--(-2.33992,1.16996);
\draw[->,structurecolor] (-2,2)--(-2.28732,2.28732);
\draw[->,structurecolor] (-1,-2)--(-1.16996,-2.33992);
\draw[->,structurecolor] (-1,-1)--(-1.22782,-1.22782);
\draw[->,structurecolor] (-1,0)--(-1.27500,0.00000);
\draw[->,structurecolor] (-1,1)--(-1.22782,1.22782);
\draw[->,structurecolor] (-1,2)--(-1.16996,2.33992);
\draw[->,structurecolor] (0,-2)--(0.00000,-2.36667);
\draw[->,structurecolor] (0,-1)--(0.00000,-1.27500);
\draw[->,structurecolor] (0,1)--(0.00000,1.27500);
\draw[->,structurecolor] (0,2)--(0.00000,2.36667);
\draw[->,structurecolor] (1,-2)--(1.16996,-2.33992);
\draw[->,structurecolor] (1,-1)--(1.22782,-1.22782);
\draw[->,structurecolor] (1,0)--(1.27500,0.00000);
\draw[->,structurecolor] (1,1)--(1.22782,1.22782);
\draw[->,structurecolor] (1,2)--(1.16996,2.33992);
\draw[->,structurecolor] (2,-2)--(2.28732,-2.28732);
\draw[->,structurecolor] (2,-1)--(2.33992,-1.16996);
\draw[->,structurecolor] (2,0)--(2.36667,0.00000);
\draw[->,structurecolor] (2,1)--(2.33992,1.16996);
\draw[->,structurecolor] (2,2)--(2.28732,2.28732);
\end{scope}
\begin{scope}[shift={(0,-6)}]
\draw[->,gray] (-2.3,0)--(2.4,0) node[right] {$x$};\draw[->,gray] (0,-2.3)--(0,2.4) node[right] {$y$};
\node at (0,-2.9) {$\rr/r^3$};
\draw[->,structurecolor] (-2,-2)--(-2.04321,-2.04321);
\draw[->,structurecolor] (-2,-1)--(-2.08199,-1.04099);
\draw[->,structurecolor] (-2,0)--(-2.11000,0.00000);
\draw[->,structurecolor] (-2,1)--(-2.08199,1.04099);
\draw[->,structurecolor] (-2,2)--(-2.04321,2.04321);
\draw[->,structurecolor] (-1,-2)--(-1.04099,-2.08199);
\draw[->,structurecolor] (-1,-1)--(-1.12964,-1.12964);
\draw[->,structurecolor] (-1,0)--(-1.27500,0.00000);
\draw[->,structurecolor] (-1,1)--(-1.12964,1.12964);
\draw[->,structurecolor] (-1,2)--(-1.04099,2.08199);
\draw[->,structurecolor] (0,-2)--(0.00000,-2.11000);
\draw[->,structurecolor] (0,-1)--(0.00000,-1.27500);
\draw[->,structurecolor] (0,1)--(0.00000,1.27500);
\draw[->,structurecolor] (0,2)--(0.00000,2.11000);
\draw[->,structurecolor] (1,-2)--(1.04099,-2.08199);
\draw[->,structurecolor] (1,-1)--(1.12964,-1.12964);
\draw[->,structurecolor] (1,0)--(1.27500,0.00000);
\draw[->,structurecolor] (1,1)--(1.12964,1.12964);
\draw[->,structurecolor] (1,2)--(1.04099,2.08199);
\draw[->,structurecolor] (2,-2)--(2.04321,-2.04321);
\draw[->,structurecolor] (2,-1)--(2.08199,-1.04099);
\draw[->,structurecolor] (2,0)--(2.11000,0.00000);
\draw[->,structurecolor] (2,1)--(2.08199,1.04099);
\draw[->,structurecolor] (2,2)--(2.04321,2.04321);
\end{scope}
\begin{scope}[shift={(6,-6)}]
\draw[->,gray] (-2.3,0)--(2.4,0) node[right] {$x$};\draw[->,gray] (0,-2.3)--(0,2.4) node[right] {$y$};
\node at (0,-2.9) {偶极方向};
\draw[->,structurecolor] (-2,-2)--(-1.96592,-1.98864);
\draw[->,structurecolor] (-2,-1)--(-1.94697,-1.01768);
\draw[->,structurecolor] (-2,0)--(-2.00000,-0.06111);
\draw[->,structurecolor] (-2,1)--(-2.05303,0.98232);
\draw[->,structurecolor] (-2,2)--(-2.03408,2.01136);
\draw[->,structurecolor] (-1,-2)--(-0.94933,-1.94088);
\draw[->,structurecolor] (-1,-1)--(-0.81291,-0.93764);
\draw[->,structurecolor] (-1,0)--(-1.00000,-0.27500);
\draw[->,structurecolor] (-1,1)--(-1.18709,1.06236);
\draw[->,structurecolor] (-1,2)--(-1.05067,2.05912);
\draw[->,structurecolor] (0,-2)--(0.00000,-1.89000);
\draw[->,structurecolor] (0,-1)--(0.00000,-0.63333);
\draw[->,structurecolor] (0,1)--(0.00000,1.36667);
\draw[->,structurecolor] (0,2)--(0.00000,2.11000);
\draw[->,structurecolor] (1,-2)--(0.94933,-1.94088);
\draw[->,structurecolor] (1,-1)--(0.81291,-0.93764);
\draw[->,structurecolor] (1,0)--(1.00000,-0.27500);
\draw[->,structurecolor] (1,1)--(1.18709,1.06236);
\draw[->,structurecolor] (1,2)--(1.05067,2.05912);
\draw[->,structurecolor] (2,-2)--(1.96592,-1.98864);
\draw[->,structurecolor] (2,-1)--(1.94697,-1.01768);
\draw[->,structurecolor] (2,0)--(2.00000,-0.06111);
\draw[->,structurecolor] (2,1)--(2.05303,0.98232);
\draw[->,structurecolor] (2,2)--(2.03408,2.01136);
\end{scope}
\end{tikzpicture}\end{center}

图中箭头在各采样点从该点起指向场的方向, 为避免奇点附近箭头挤叠, 以单调尺度$\ell_{\rm arrow}=0.55|\vct f|/(1+|\vct f|)$压缩长度. 不把这种绘图长度当真实场强; 原点有奇点的两场未画该点. 常矢量向右下, 径向场在不同半径的增长/衰减和偶极轴外转向都可直接看见.

平面源场$\vct v_s=a(x\uv x+y\uv y)$有散度$2a$、旋度零; 环流场$\vct v_c=\Omega(-y\uv x+x\uv y)$有散度零、旋度$2\Omega\uv z$. 二者相加可同时有源与旋转. 下图取$a=\Omega=1$只示意方向, 长度按同一比例缩短.
\begin{center}\begin{tikzpicture}[>=Stealth,scale=.72]
\foreach \X in {-1,0,1}{\foreach \Y in {-1,0,1}{
 \draw[->] (\X,\Y)--({1.3*\X},{1.3*\Y});
 \draw[->] ({\X+4},\Y)--({\X+4-.3*\Y},{\Y+.3*\X});}}
\node at (0,-1.8){源场};\node at (4,-1.8){环流场};
\end{tikzpicture}\end{center}
电场通量与磁场环量可以借用这样的语言, 但电磁场不是流体速度. 静电场在源外无旋; 磁场没有单极源而总磁通为零; 时变电场则可以有旋度. 这些物理规律要由后文的高斯与法拉第定律确定, 不是单凭流体类比证明.
\end{solution}

\originalsection{离散电荷与带电粒子}\label{tb:coulombexamples}
叠加时先明确“谁对谁作用”, 再相加分量. 同一源对观察点的位移是$\rr-\rr_i$, 力的单位N, 场的单位N/C. 三角形的几何对称与电荷符号要同时使用.
\begin{example}
\textbf{原2.1: 三电荷的合力.} 
$q_1=+6.0\,\mu\mathrm C$在$(0,0)$, $q_2=-6.0\,\mu\mathrm C$在$(0,a)$, $q_3=+3.0\,\mu\mathrm C$在$(a,a)$, $a=2.0\,\mathrm{cm}$. 求$q_3$的合力大小及其与$+x$轴的夹角.
\end{example}
\begin{solution}
$q_1$到$q_3$距离$\sqrt2a$, 单位方向$(\uv x+\uv y)/\sqrt2$; $q_2$到$q_3$距离$a$, 引力向$-x$. 所以
\[\vct F_3=\frac{\kc q_1q_3}{a^2}\left[\left(\frac{\sqrt2}{4}-1\right)\uv x+\frac{\sqrt2}{4}\uv y\right].\]
代$\kc\simeq9.0\times10^9$得$F_x=-2.62\,\mathrm N$, $F_y=1.43\,\mathrm N$, 大小约$3.0\,\mathrm N$, 方向$151.3^\circ$. 普通反正切给$-28.7^\circ$还需判象限; 第二象限由$F_x<0,F_y>0$决定. 不应把不同方向的两个带符号力大小直接相加.
\end{solution}
\begin{example}\label{ex:hydrogen}
原2.13.1给经典氢原子半径$r=0.53\times10^{-10}$m, 问电子与质子的电力、质子的电场、电力与引力之比, 及为何行星运动通常由引力支配.
\end{example}
\begin{solution}
按原题近似$e=1.6\times10^{-19}$C、$\kc=9.0\times10^9$N$\cdot$m$^2$/C$^2$, $F_e=\kc e^2/r^2\simeq8.2\times10^{-8}$N, $E=\kc e/r^2\simeq5.1\times10^{11}$N/C. 原答案用另一舍入半径写出约$5.76\times10^{11}$N/C; 本稿保持题面的$0.53\times10^{-10}$m一致计算. 比值
$F_e/F_g=\kc e^2/(Gm_em_p)\simeq2.2\times10^{39}$, $r$约去. 宏观物体近乎电中性, 正负电作用大量抵消, 质量引力却同向相加. 这不表示带电宏观物体没有重要电力. 经典电子绕核运动也不是稳定原子的完整模型; 本例只比较两种力的尺度.
\end{solution}

\begin{example}\label{ex:millikan}
原2.13.2给球形油滴$r=1.64\times10^{-6}$m, 密度$8.51\times10^2$kg/m$^3$, 向下电场$1.92\times10^5$N/C恰好平衡重力. 忽略空气浮力, 求质量、带电量与电子数.
\end{example}
\begin{solution}
质量$M=4\pi\rho_{\rm oil}r^3/3\simeq1.57\times10^{-14}$kg. 电力必须向上而场向下, 故$q<0$. 平衡给$q=-Mg/E\simeq-8.03\times10^{-19}$C, 约为$-5e$. 原题把$e$作为正的基本电荷大小, 所以是约五个多余电子, 不能写成油滴正带电. 真正实验还需计浮力、黏滞和测量误差; 单个例子的接近整数不能证明任意电荷量子化.
\end{solution}

\begin{example}\label{ex:electronplates}
原2.13.3: 电子以水平速度$v_0$进入向下均匀场$E$, 极板长$L_1$, 出板后再水平前进$L_2$到屏. 求力、加速度、出板时间、速度、位移、偏角和总偏移. 假定低速、忽略重力及边缘场, 且电子没有先撞上极板.
\end{example}
\begin{solution}
取$+y$向上, $F_y=eE$, $a_y=eE/m_e$, $x=v_0t$, 所以$t_1=L_1/v_0$.
出板时$v_x=v_0$, $v_y=eEL_1/(m_ev_0)$, $y_1=eEL_1^2/(2m_ev_0^2)$,
$\tan\theta_1=eEL_1/(m_ev_0^2)$. 出板后的匀速时间为$L_2/v_0$, 因而
\[
y_{\rm screen}=\frac{eEL_1}{m_ev_0^2}\left(\frac{L_1}{2}+L_2\right).
\]
这是两段运动的接续: 第一段有恒加速度, 第二段保留竖直速度. 若总位移超出板间允许空间, 原路径假设失效, 应先求碰板时间.
\end{solution}
